% Manuscript-ready module. Requires amsmath, amssymb, theorem environments.
% Source status: deductions beyond the two pinned incoming reports;
% classical Faulhaber, Hurwitz, Lerch and polynomial integration inputs
% are explicitly identified. No literature-wide novelty claim is made.

\section{Nonpositive harmonic indices and a convergent Lerch generator}
\label{harm:section}

The incoming independent-order report reduces a depth-two word whose
last index is a nonpositive integer. We extend that calculation to an
arbitrary position and depth, prove a canonical finite reduction, and
sum all zero-index decorations around one free index in a convergent
Lerch generating function. The summation of powers by Bernoulli
polynomials is classical; the contribution here is its compatibility
with the entire relative harmonic interpolant at arbitrary endpoints,
including its independent spectral derivatives. The interpolant itself
is established in the incoming report, and its one-endpoint version is
the interpolant of Ihara, Nakamura, and Yamamoto
\cite{harm:INY}. We do not claim literature-wide priority for these
consequences.

\subsection{Normalization and a local elimination theorem}

For $\Re a,\Re b>0$, use the normalization
\begin{align}
 Z_a(s_1,\ldots,s_r)
  &=\sum_{n_1>\cdots>n_r\ge0}
                  \prod_{j=1}^r(n_j+a)^{-s_j},\notag\\
 \mathcal H_{a,b}(w)
  &=Z_b(w)-\sum_{j=1}^{|w|}Z_a(w_{\le j})
                              \mathcal H_{a,b}(w_{>j}),
 \qquad \mathcal H_{a,b}(\varnothing)=1.
 \label{harm:normalization}
\end{align}
The $Z_a$ are meromorphically continued from their convergence domain.
The combined function $\mathcal H_{a,b}$ is entire in all spectral
indices. For $a=b+N$, $N\ge0$ an integer, it is the strict finite sum
over $N>n_1>\cdots>n_r\ge0$. In particular,
\begin{equation}
 h_{a,b}(s):=\mathcal H_{a,b}(s)
      =\zeta(s,b)-\zeta(s,a),\qquad
 h_{a,b}(1)=\psi(a)-\psi(b).
 \label{harm:h}
\end{equation}
We take $B_n=B_n(0)$, so $B_1=-1/2$, whereas $B_1(1)=1/2$.
This distinction is essential for strict inequalities.

\begin{theorem}[Eliminating a fixed nonpositive index]
\label{harm:local}
Let $m\ge0$ and $q=m+1$. All other indices below are arbitrary
independent complex numbers. With unchanged prefix and suffix words
suppressed, an interior index obeys
\begin{align}
 \mathcal H_{a,b}(\ldots,s,-m,t,\ldots)
  =\frac1q\sum_{k=0}^q\binom qk\bigl\{
   &B_{q-k}\mathcal H_{a,b}(\ldots,s-k,t,\ldots)\notag\\[-2pt]
   &-B_{q-k}(1)\mathcal H_{a,b}(\ldots,s,t-k,\ldots)
   \bigr\}.
 \label{harm:interior}
\end{align}
At the two boundaries the formulas are
\begin{align}
 \mathcal H_{a,b}(-m,t,\ldots)
 &=\frac1q\left\{B_q(a)\mathcal H_{a,b}(t,\ldots)
  -\sum_{k=0}^q\binom qk B_{q-k}(1)
                        \mathcal H_{a,b}(t-k,\ldots)\right\},
 \label{harm:top}\\
 \mathcal H_{a,b}(\ldots,s,-m)
 &=\frac1q\left\{\sum_{k=0}^q\binom qk B_{q-k}
                  \mathcal H_{a,b}(\ldots,s-k)
                     -B_q(b)\mathcal H_{a,b}(\ldots,s)\right\}.
 \label{harm:bottom}
\end{align}
The singleton rule is
\begin{equation}
             h_{a,b}(-m)=\frac{B_q(a)-B_q(b)}q.
 \label{harm:singleton}
\end{equation}
These are identities of entire functions in every remaining index.
\end{theorem}

\begin{proof}
The classical finite power-sum identity \cite{DLMF} gives
\begin{equation}
 \sum_{z<x<y}x^m=\frac{B_q(y)-B_q(z+1)}q,
 \label{harm:finite-gap}
\end{equation}
when the variables belong to a common unit lattice and the sum has the
indicated strict endpoints. The upper and lower boundary variants
replace $y$ by $a$ and $z+1$ by $b$, respectively. This explains the
formulas for finite sums, but agreement at integer endpoint differences
alone would not justify their analytic continuation. We instead prove
the identities directly in \eqref{harm:normalization}.

For $Z_a$, summing an interior or final nonpositive index by
\eqref{harm:finite-gap} proves \eqref{harm:interior} and
\eqref{harm:bottom}, with $\mathcal H_{a,b}$ replaced by $Z_a$ and
the final lower endpoint $b$ replaced by $a$. One first chooses the
unfixed indices in a domain of absolute convergence and then continues
meromorphically. For a first index use
\[
 \sum_{n_1>n_2}(n_1+a)^{-z}
                   =\zeta(z,n_2+a+1),\qquad
 \zeta(-m,x+1)=-\frac{B_q(x+1)}q.
\]
Taking the remaining indices sufficiently far into their convergence
domain justifies continuation of $z$ to $-m$ inside this expression.
Consequently the top-index identity for $Z_a$ is
\begin{equation}
 Z_a(-m,t,\ldots)
   =-\frac1q\sum_{k=0}^q\binom qk B_{q-k}(1)
                               Z_a(t-k,\ldots).
 \label{harm:Ztop}
\end{equation}
All these statements concern generic values on the indicated spectral
hyperplane, followed by meromorphic continuation; they assign no
directional value to an individual singular $Z_a$ at an intersection.

Now induct on the depth in \eqref{harm:normalization}. For the top
formula substitute \eqref{harm:Ztop} into every prefix of length at
least two. The prefix of length one contributes
$-Z_a(-m)\mathcal H_{a,b}(t,\ldots)
 =B_q(a)\mathcal H_{a,b}(t,\ldots)/q$.
Each remaining group is exactly a shorter convolution identity, giving
\eqref{harm:top}. The bottom formula follows in the same way. Its
two exceptional terms are the last prefix and the singleton suffix;
the factors involving $B_q(a)$ cancel, leaving $B_q(b)$.

For clarity, consider an interior letter and write its nonempty words
on the two sides as $L$ and $R$. In the deconcatenation sum, cuts
strictly before $L$'s last letter use the induction hypothesis on the
suffix; cuts after the first letter of $R$ use the interior $Z_a$
identity on the prefix. The two adjacent cuts are
\[
 Z_a(L)\mathcal H_{a,b}(-m,R),\qquad
 Z_a(L,-m)\mathcal H_{a,b}(R).
\]
After the top and bottom substitutions their exceptional terms are
$B_q(a)Z_a(L)\mathcal H_{a,b}(R)/q$ and its negative.
They cancel. The remaining left-shifted and right-shifted terms are
precisely the recursions for the two shorter words in
\eqref{harm:interior}. This proves that formula on generic remaining
indices. Entireness of $\mathcal H_{a,b}$ then extends all the
identities to every remaining index, including coincident resonances.
\end{proof}

For $m=0$, the interior rule takes the particularly transparent form
\begin{equation}
 \mathcal H_{a,b}(\ldots,s,0,t,\ldots)
 =\mathcal H_{a,b}(\ldots,s-1,t,\ldots)
  -\mathcal H_{a,b}(\ldots,s,t-1,\ldots)
  -\mathcal H_{a,b}(\ldots,s,t,\ldots).
 \label{harm:zero}
\end{equation}
The last term records the strict gap's missing endpoint. Changing
$B_1(1)$ to $B_1$ incorrectly deletes it.

\subsection{A canonical polynomial normal form}

Suppose a word contains $k$ free indices, in their original order
$s_1,\ldots,s_k$, and fixed indices $-m_1,\ldots,-m_e$, where
$m_j\ge0$. Set $M=\sum_{j=1}^e(m_j+1)$.

\Needspace{13\baselineskip}
\begin{theorem}[Finite reduction at arbitrary depth]
\label{harm:normalform}
There are explicitly computable polynomials $C_{\boldsymbol\nu}(a,b)
\in\mathbb Q[a,b]$ such that
\begin{equation}
 \boxed{\displaystyle
 \mathcal H_{a,b}(w)=
  \sum_{\substack{\boldsymbol\nu\in\mathbb N_0^k\\
                    |\boldsymbol\nu|\le M}}
  C_{\boldsymbol\nu}(a,b)
  \mathcal H_{a,b}(s_1-\nu_1,\ldots,s_k-\nu_k),\qquad
 \deg C_{\boldsymbol\nu}\le M-|\boldsymbol\nu|.}
 \label{harm:normalform-eq}
\end{equation}
For $k=0$ the result is a polynomial multiplying the empty word.
Every elimination order in Theorem~\ref{harm:local} gives the same
coefficient polynomials when the free indices are kept formally
distinct. The bound concerns this explicitly specified finite normal
form and does not assert a minimal arithmetic depth.
\end{theorem}

\begin{proof}
Every local rule deletes one fixed index. A deleted index $-m$ has
budget $q=m+1$: it either produces an endpoint polynomial of degree
at most $q$, or shifts an adjacent index by $k\le q$. If the adjacent
index is also fixed, its later budget increases by $k$; if it is free,
its final shift increases by $k$. Induction therefore proves termination
and the degree bound in \eqref{harm:normalform-eq}.

Here is a canonical description of the coefficients. For a fixed
nonpositive block $v$, let
\[
 \Pi_v(A,B)=\mathcal H_{A,B}(v),\qquad \Pi_{\varnothing}(A,B)=1.
\]
The first part proves that these are polynomials. Write
$w=(v_0,s_1,v_1,\ldots,s_k,v_k)$ and expand the polynomial
\begin{equation}
 G_w(\boldsymbol y;a,b)=
   \Pi_{v_0}(a,y_1+1)
   \prod_{j=1}^{k-1}\Pi_{v_j}(y_j,y_{j+1}+1)
   \Pi_{v_k}(y_k,b)
   =\sum_{\boldsymbol\nu}C_{\boldsymbol\nu}(a,b)
                                 \boldsymbol y^{\boldsymbol\nu}.
 \label{harm:gaps}
\end{equation}
For integer endpoint differences, fix the free summation variables
$y_1>\cdots>y_k$. The sums in the intervening gaps are independent,
and their product is exactly \eqref{harm:gaps}.

To identify these coefficients with those of any already proved local
reduction, compare the two finite sums for all independent $s_j$.
Distinct ordered tuples of positive lattice points give linearly
independent functions $\prod_j y_j^{-s_j}$: this is elementary finite
linear independence of exponential functions, applied successively
in the $s_j$. Thus the difference of the two coefficient polynomials
vanishes at every lattice configuration
$a=b+N$, $y_j=b+n_j$, $N>n_1>\cdots>n_k\ge0$, with $b>0$.
These configurations are Zariski dense. Indeed a polynomial vanishing
on all integer tuples satisfying strict inequalities vanishes by
successive univariate polynomial uniqueness, and $b$ ranges over an
interval. Hence the coefficient polynomials agree identically.
This is polynomial uniqueness after the analytic identities have
already been proved; it is not uniqueness of a holomorphic interpolation
from integer endpoint values.
\end{proof}

In the one-free-index case, every depth reduces to the finite sum
\begin{equation}
           \mathcal H_{a,b}(v_0,s,v_1)
              =\sum_{\nu=0}^M C_\nu(a,b)h_{a,b}(s-\nu).
 \label{harm:onefree}
\end{equation}
All derivatives in the free index commute with this identity. At
$s-\nu=1$ use the finite value
\begin{equation}
 h_{a,b}^{(r)}(1)=(-1)^r\{\gamma_r(b)-\gamma_r(a)\},
 \qquad
 h_{a,b}'(0)=\log\Gamma(b)-\log\Gamma(a).
 \label{harm:jets}
\end{equation}
At positive integer arguments $>1$, the undifferentiated differences
are polygamma differences. These standard Hurwitz specializations are
recorded in \cite{DLMF}. Differentiation of a \emph{fixed}
index $-m$ is a different problem and is not licensed by a formula
proved only on that spectral hyperplane.

Two concise examples are
\begin{align}
 \mathcal H_{a,b}(0,s,0)
   &=-h_{a,b}(s-2)+(a+b-1)h_{a,b}(s-1)
                                  -b(a-1)h_{a,b}(s),
 \label{harm:examplezero}\\
 \mathcal H_{a,b}(-1,s,-1)
   &=\frac14\bigl\{-h_{a,b}(s-4)
       +(A+B+1)h_{a,b}(s-2)\notag\\
   &\hspace{32mm}+(B-A)h_{a,b}(s-1)-ABh_{a,b}(s)\bigr\}.
 \label{harm:exampleone}
\end{align}
Here $A=a(a-1)$ and $B=b(b-1)$ in the second identity.
The absent shift $s-3$ in the second identity is an exact cancellation.
It follows by expanding
$(A-y^2-y)(y^2-y-B)/4$. For example, its derivative at $s=2$ is
\begin{align}
 \left.\partial_s\mathcal H_{a,b}(-1,s,-1)\right|_{s=2}
  =\frac14\bigl\{&-\zeta'(-2,b)+\zeta'(-2,a)\notag\\
       &+(A+B+1)\log\tfrac{\Gamma(b)}{\Gamma(a)}
        +(B-A)\{\gamma_1(a)-\gamma_1(b)\}\notag\\
       &-AB\{\zeta'(2,b)-\zeta'(2,a)\}\bigr\}.
 \label{harm:examplejet}
\end{align}

\subsection{Summing arbitrarily many zero indices}

For a block of $p$ zero indices, polynomial continuation of the finite
count gives
\begin{equation}
                    \Pi_{\{0\}^p}(A,B)=\binom{A-B}{p}.
 \label{harm:zeroblock}
\end{equation}
Consequently \eqref{harm:gaps} for a single free index uses the polynomial
$\binom{a-y-1}{p}\binom{y-b}{q}$. All such polynomials can be summed
in one analytic formula.

\begin{theorem}[Lerch generator for zero decorations]
\label{harm:Lerch}
Let $\Re a,\Re b>0$. There is a function
$\mathcal L_{a,b}(s;t)$ entire in $s$ and holomorphic for $|t|<2\pi$
such that, when $\Re t<0$,
\begin{equation}
 \mathcal L_{a,b}(s;t)=
    e^{bt}\Phi(e^t,s,b)-e^{at}\Phi(e^t,s,a).
 \label{harm:Lerchdifference}
\end{equation}
Here $\Phi(z,s,c)=\sum_{n\ge0}z^n(n+c)^{-s}$ for $|z|<1$, with
principal powers of $n+c$. The combination in
\eqref{harm:Lerchdifference}, not its separate summands, is continued
through $t=0$. It satisfies
\begin{equation}
 \mathcal L_{a,b}(s;0)=h_{a,b}(s),\qquad
 \partial_t\mathcal L_{a,b}(s;t)=\mathcal L_{a,b}(s-1;t),\qquad
 \mathcal L_{a,b}(s;t)=\sum_{n\ge0}h_{a,b}(s-n)\frac{t^n}{n!}.
 \label{harm:LerchTaylor}
\end{equation}
The last series converges locally normally for $|t|<2\pi$, including
all fixed spectral derivatives. For $|u|,|v|<1/2$ one has the convergent
identity
\begin{equation}
 \boxed{\displaystyle
 \sum_{p,q\ge0}\mathcal H_{a,b}(\{0\}^p,s,\{0\}^q)u^pv^q
  =(1+u)^{a-1}(1+v)^{-b}
    \mathcal L_{a,b}\bigl(s;\log(1+v)-\log(1+u)\bigr).}
 \label{harm:zerogenerator}
\end{equation}
Powers and logarithms on the right are their germs at $u=v=0$.
\end{theorem}

\begin{proof}
Define
\begin{equation}
 K_{a,b}(y)=\frac{e^{-by}-e^{-ay}}{1-e^{-y}}.
 \label{harm:K}
\end{equation}
The apparent pole at $y=0$ is removable, with value $a-b$.
Its only other possible finite poles are $2\pi i\mathbb Z\setminus\{0\}$.
Thus for $|t|<2\pi$ the function $K_{a,b}(x-t)$ is analytic near
every $x\ge0$ and decreases exponentially as $x\to+\infty$,
locally uniformly in $t,a,b$. Initially for $\Re s>0$ put
\begin{equation}
 \mathcal L_{a,b}(s;t)=\frac1{\Gamma(s)}
                \int_0^\infty x^{s-1}K_{a,b}(x-t)\,dx.
 \label{harm:LerchMellin}
\end{equation}
For any $N\ge1$, subtract the Taylor polynomial of $K_{a,b}(x-t)$
at $x=0$ in the integral from $0$ to $1$, and add
\[
 \frac1{\Gamma(s)}\sum_{j=0}^{N-1}
             \frac{K_{a,b}^{(j)}(-t)}{j!(s+j)}.
\]
The subtracted integral is holomorphic for $\Re s>-N$. The displayed
possible poles are cancelled by the simple zeros of $1/\Gamma(s)$.
Since $N$ is arbitrary, this proves joint entireness in $s$ and
holomorphy for $|t|<2\pi$, with locally normal spectral derivatives.

For $\Re t<0$, geometric expansion in \eqref{harm:LerchMellin}
and the Gamma integral give \eqref{harm:Lerchdifference}, first for
$\Re s>0$ and then for all $s$. At $t=0$ the ordinary Hurwitz
integral gives $h_{a,b}(s)$, initially for $\Re s>1$, and then by
continuation. Differentiating the absolutely convergent Lerch series
for $\Re t<0$ proves the order-lowering identity in
\eqref{harm:LerchTaylor}; continuation and Taylor's theorem give the
remaining assertions.

The right side of \eqref{harm:zerogenerator} is holomorphic for
$|u|,|v|<1/2$, because the absolute value of the difference of its
logarithms is less than $2\log2<2\pi$. To identify a coefficient,
formally replace $h_{a,b}(s-n)$ in \eqref{harm:LerchTaylor} by
$y^n$. Its right side becomes $e^{yt}$. Including the prefactors in
\eqref{harm:zerogenerator} gives
\[
                 (1+u)^{a-y-1}(1+v)^{y-b}.
\]
Its coefficient of $u^pv^q$ is precisely
$\binom{a-y-1}{p}\binom{y-b}{q}$. At each fixed bidegree this is a
finite polynomial calculation. The canonical gap formula then identifies
it with $\mathcal H_{a,b}(\{0\}^p,s,\{0\}^q)$, proving the
identity and the convergence of its left side.
\end{proof}

The branch cancellation can also be read as a Gamma separation:
$\Gamma(s)\mathcal L_{a,b}(s;t)$ is the ordinary Mellin transform
of the explicit kernel \eqref{harm:K}; its apparent spectral poles at
nonpositive integers are entirely accounted for by $\Gamma(s)$.
This proves a generating formula for the specified one-free-order
family. It does not settle a minimality question for several arbitrary
independent nonzero indices.

Setting $u=v$ gives the exact insertion identity
\begin{equation}
 \sum_{p+q=r}\mathcal H_{a,b}(\{0\}^p,s,\{0\}^q)
                  =\binom{a-b-1}{r}h_{a,b}(s).
 \label{harm:diagonalzero}
\end{equation}
The generating function also has normalized antiderivatives:
\begin{equation}
 \int_0^t\mathcal L_{a,b}(s;\tau)\,d\tau
            =\mathcal L_{a,b}(s+1;t)-h_{a,b}(s+1).
 \label{harm:Lerchprimitive}
\end{equation}
At a nonpositive order,
\begin{equation}
 \mathcal L_{a,b}(-m;t)
      =\partial_t^m\frac{e^{bt}-e^{at}}{1-e^t},
 \qquad m\ge0,
 \label{harm:Lerchnegative}
\end{equation}
where $t=0$ is removable. At $a=b+N$, $N\ge0$ integer,
$\mathcal L_{a,b}(s;t)=\sum_{n=0}^{N-1}(n+b)^{-s}e^{(n+b)t}$;
the double generating function becomes a polynomial of total degree at
most $N-1$ when $N\ge1$ and is zero when $N=0$.

\subsection{Polynomial endpoint primitives and Stieltjes moments}

The polynomial Hurwitz integration mechanism used here is classical:
Espinosa and Moll give a general integration-by-parts formula and its
polynomial specialization \cite[Theorem~2.1 and Example~2.2]{harm:EM}.
The canonical manuscript \cite{source:manuscript}, in the Moments section
of \src{chapters/07-integration.tex}, already states the unit-interval
Hurwitz transform, gives the first Stieltjes moment, and identifies the complete
homogeneous coefficients for every Stieltjes order. We state the fully
indexed arbitrary-endpoint version, with the normalization required by
\eqref{harm:onefree}. This is an explicit extension and restatement of
the existing moment mechanism, and is not counted as a new general
Stieltjes-moment principle.

For $d\ge0$, write $\Delta=a-b$ and define
\begin{equation}
 P_{d;a,b}(s)=\int_b^a(x-b)^d h_{x,b}(s)\,dx.
 \label{harm:Pdef}
\end{equation}
Paths lie in the right half-plane. The integral is path independent and
entire in $s$. Initially away from singularities of the individual
displayed terms,
\begin{align}
 P_{d;a,b}(s)
  ={}&\frac{\Delta^{d+1}}{d+1}\zeta(s,b)\notag\\
   &+\sum_{j=0}^d\frac{d!}{(d-j)!}
     \frac{\Delta^{d-j}\zeta(s-j-1,a)
              -\mathbf1_{j=d}\zeta(s-j-1,b)}
          {(s-1)(s-2)\cdots(s-j-1)}.
 \label{harm:Pfinite}
\end{align}
This follows by repeated integration by parts using
$\partial_x\zeta(s-1,x)=-(s-1)\zeta(s,x)$.
The definite integral proves that every apparent pole in the combined
right side is removable. Expanding a coefficient polynomial
$C_\nu(x,b)$ in powers of $x-b$ and using \eqref{harm:Pfinite}
therefore integrates every one-free-order word and every one of its
free-order derivatives in a finite Hurwitz basis, with the constant
fixed at $a=b$.

To state all resonance coefficients, put
\begin{equation}
 c_{j,n}=[z^n]\prod_{\ell=1}^j(1-z/\ell)^{-1},
 \qquad c_{0,n}=\mathbf1_{n=0}.
 \label{harm:cn}
\end{equation}
These are finite generalized harmonic expressions:
\[
 \sum_{n\ge0}c_{j,n}z^n
     =\exp\left(\sum_{r\ge1}\frac{H_j^{(r)}}r z^r\right),
 \qquad H_j^{(r)}=\sum_{\ell=1}^j\ell^{-r}.
\]
For example $c_{j,0}=1$, $c_{j,1}=H_j$ and
$c_{j,2}=(H_j^2+H_j^{(2)})/2$. Equivalently, $c_{j,n}$ is the
complete exponential Bell polynomial in
$H_j,1!H_j^{(2)},\ldots,(n-1)!H_j^{(n)}$, divided by $n!$.

\begin{corollary}[Explicit arbitrary-endpoint form of the moment calculus]
\label{harm:Stieltjesmoments}
For $d,m\ge0$ and $\Re a,\Re b>0$,
\begin{align}
 \int_b^a(x-b)^d\gamma_m(x)\,dx
   ={}&(-1)^{m+1}m!
        \sum_{j=0}^d(-1)^j\binom dj
        \sum_{r=0}^{m+1}\frac{c_{j,m+1-r}}{r!}\notag\\
      &\hspace{7mm}\cdot
       \left\{\Delta^{d-j}\zeta^{(r)}(-j,a)
                      -\mathbf1_{j=d}\zeta^{(r)}(-j,b)\right\}.
 \label{harm:Stieltjesfinite}
\end{align}
Here derivatives of $\zeta$ are with respect to its first argument.
\end{corollary}

\begin{proof}
Set $s=1+z$ in \eqref{harm:Pfinite}. The denominator in its $j$th
term has expansion
\[
 \frac1{z(z-1)\cdots(z-j)}
       =\frac{(-1)^j}{j!z}\sum_{n\ge0}c_{j,n}z^n.
\]
Use
$\zeta(1+z,b)=z^{-1}+\sum_{m\ge0}(-1)^m\gamma_m(b)z^m/m!$
and the ordinary Taylor expansion of each $\zeta(z-j,a)$.
The coefficient of $z^{-1}$ cancels, as it must by entireness of
$P_{d;a,b}$. The coefficient of $z^m$, multiplied by $m!$, equals
\[
 (-1)^m\frac{\Delta^{d+1}}{d+1}\gamma_m(b)
 +m!\sum_{j=0}^d(-1)^j\binom dj
    \sum_{r=0}^{m+1}\frac{c_{j,m+1-r}}{r!}
       \left\{\Delta^{d-j}\zeta^{(r)}(-j,a)
                         -\mathbf1_{j=d}\zeta^{(r)}(-j,b)\right\}.
\]
On the integral side, \eqref{harm:jets} gives
$(-1)^m\{\Delta^{d+1}\gamma_m(b)/(d+1)
 -\int_b^a(x-b)^d\gamma_m(x)\,dx\}$.
Cancel the common endpoint term and solve for the integral.
All coefficient extractions are locally normal on compact endpoint
sets in the right half-plane.
\end{proof}

For $d=0$ this recovers the familiar normalized primitive
\begin{equation}
 \int_b^a\gamma_m(x)\,dx
       =\frac{(-1)^{m+1}}{m+1}
          \{\zeta^{(m+1)}(0,a)-\zeta^{(m+1)}(0,b)\}.
 \label{harm:d0}
\end{equation}
For $d=1$, since $c_{1,n}=1$ for every $n$,
\begin{align}
 \int_b^a(x-b)\gamma_m(x)\,dx
  =(-1)^{m+1}m!\left\{
       \frac{\Delta\zeta^{(m+1)}(0,a)}{(m+1)!}
       -\sum_{r=0}^{m+1}
          \frac{\zeta^{(r)}(-1,a)-\zeta^{(r)}(-1,b)}{r!}
                         \right\}.
 \label{harm:d1}
\end{align}
In particular $m=0$ reduces to an integral of the digamma function,
and Lerch's identity converts the $\zeta'(0,a)$ term to
$\log\Gamma(a)-\tfrac12\log(2\pi)$.

\subsection{Exact scope and further questions}

The identities above add arbitrary-depth finite reductions and a
convergent all-zero-decoration generator to the previously established
depth-two calculus. They do not evaluate nonsymmetric positive-index
coordinates such as the incoming harmonic Stieltjes parameter $\eta$.
No statement about the unresolved $S_6$ or $S_8$ reductions is made.

Natural next questions are the following.
\begin{enumerate}
\item Construct and prove a jointly convergent analogue of
\eqref{harm:zerogenerator} with two or more independently moving free
indices and arbitrary zero blocks between them. The finite coefficient
polynomials \eqref{harm:gaps} are already explicit; the remaining task
is a useful analytic kernel and a precise simultaneous branch domain.
\item Differentiate an eliminated index before specializing it to
$-m$. Determine a finite set of additional ordered coordinates needed
at the first such derivative. The present identities only differentiate
the free indices, so this question asks for a real extension.
\item For rational endpoints, combine \eqref{harm:Stieltjesfinite}
with finite character projections to give algorithms for polynomial
moments in Dirichlet $L$-derivatives, retaining every conductor logarithm.
\item Formalize the Bernoulli deletion rules, the polynomial gap
normal form, and the finite coefficient extraction. Their analytic
interface is the already established entire interpolant and the
explicit normalized Mellin continuation in \eqref{harm:LerchMellin}.
\end{enumerate}
